\subsection{Rancangan Arsitektur Solusi}
\label{subsec:rancangan-arsitektur-solusi}

Arsitektur solusi mempertahankan tiga \textit{client} utama DecMed, yaitu \textit{Patient Client}, \textit{Hospital Client}, dan \textit{Ministry Client}. \textit{Patient Client} digunakan pasien untuk menerbitkan token akses awal. \textit{Hospital Client} digunakan oleh petugas administrasi, dokter, dan petugas laboratorium untuk mengakses atau memperbarui data sesuai kewenangannya. \textit{Ministry Client} tetap berada pada ruang lingkup manajemen fasyankes sebagaimana rancangan DecMed awal.

Perubahan utama berada pada \textit{server} PRE dan struktur token akses. Pada DecMed awal, token akses masih membawa informasi yang relatif terbatas, seperti \textit{address}, \textit{role}, dan \textit{purpose}. Pada rancangan ini, token tersebut diganti dengan Macaroon yang membawa daftar \textit{caveats}. Setiap \textit{request} akses data harus menyertakan Macaroon. \textit{Server} PRE kemudian memvalidasi keaslian token, mengevaluasi seluruh \textit{caveats}, memastikan bahwa permintaan akses sesuai dengan segmen data yang diminta, lalu menjalankan proses \textit{re-encryption} hanya jika semua pemeriksaan berhasil.

Gambaran arsitektur sistem secara keseluruhan dapat dilihat pada gambar X
% TODO: Gambar arsitektur solusi
, sedangkan rancangan penyimpanan \textit{on-chain} dan \textit{off-chain} dijelaskan pada Tabel X
dan Tabel Y
% TODO: Tabel skema penyimpanan on-chain dan off-chain